\chapter{Authorization: Roles and Permissions}

Knowing \emph{who} made a request is only half the story -- knowing \emph{what they're allowed to do} is the other half. A user and an admin both pass \texttt{protect} from Chapter 20, but only one should delete someone else's task.

\section{Authentication vs Authorization}

\begin{center}
\begin{tabularx}{\textwidth}{lXX}
\toprule
\textbf{} & \textbf{Authentication} & \textbf{Authorization} \\
\midrule
Question & "Who are you?" & "What can you do?" \\
When & Every protected request, first & After authentication \\
Implemented by & \texttt{protect} (Ch. 18/20) & \texttt{authorize(...roles)} (this chapter) \\
Failure & 401 Unauthorized & 403 Forbidden \\
Based on & Valid JWT/session & \texttt{req.user.role} \\
\bottomrule
\end{tabularx}
\end{center}

\begin{notebox}[NOTE]
Shortcut: \textbf{401} means "I don't know who you are." \textbf{403} means "I know who you are, and the answer is no."
\end{notebox}

\section{The Role Field}

Recall the User schema's \texttt{role} field (Chapter 13):

\begin{lstlisting}[language=JavaScript]
// models/User.js (relevant field)
role: { type: String, enum: ["user", "manager", "admin"], default: "user" },
\end{lstlisting}

\begin{itemize}
  \item \textbf{user} -- own tasks only.
  \item \textbf{manager} -- team tasks.
  \item \textbf{admin} -- full access.
\end{itemize}

\section{The \texttt{authorize} Middleware Factory}

One factory function takes allowed roles and returns a tailored middleware, instead of a separate middleware per role combination.

\begin{lstlisting}[language=JavaScript]
// middleware/authorize.js
export const authorize = (...allowedRoles) => {
  return (req, res, next) => {
    if (!req.user) {
      return res.status(401).json({ success: false, message: "Not authenticated" });
    }

    if (!allowedRoles.includes(req.user.role)) {
      return res.status(403).json({
        success: false,
        message: `Role '${req.user.role}' is not permitted to perform this action`,
      });
    }

    next();
  };
};
\end{lstlisting}

\texttt{authorize} always runs after \texttt{protect}, so \texttt{!req.user} is just a defensive guard.

\section{Using It on Routes}

\begin{lstlisting}[language=JavaScript]
// routes/taskRoutes.js
import { protect } from "../middleware/authMiddleware.js";
import { authorize } from "../middleware/authorize.js";
import { deleteTask, getTasks } from "../controllers/taskController.js";

router.get("/tasks", protect, getTasks);

router.delete("/tasks/:id", protect, authorize("admin", "manager"), deleteTask);
\end{lstlisting}

A request must pass \texttt{protect} then \texttt{authorize("admin", "manager")} before \texttt{deleteTask} runs -- a plain \texttt{user} gets a 403 and the delete never touches the database.

\begin{tipbox}[TIP]
Because roles are arguments, the same middleware covers every combination: \texttt{authorize("admin")}, \texttt{authorize("admin", "manager")}, etc.
\end{tipbox}

\begin{warnbox}[WARNING]
Role checks alone don't cover ownership -- a manager passing \texttt{authorize("manager")} shouldn't necessarily delete another team's task. Add an explicit ownership check inside the controller too.
\end{warnbox}

\whatwhyhow
{An \texttt{authorize(...roles)} middleware factory that checks \texttt{req.user.role} against a list of permitted roles for a given route.}
{Authentication only proves identity; without a separate authorization check, any logged-in user -- regardless of role -- could call sensitive endpoints like deleting tasks or managing other users.}
{Chain it after \texttt{protect} in the route definition: \texttt{router.delete("/tasks/:id", protect, authorize("admin", "manager"), deleteTask)}, returning 403 for any role not in the allowed list.}
